% id: 114-06
% book: 베이직쎈 확률과 통계 · 06 이항분포와 정규분포
% section: 유형 067 정규분포곡선의 성질을 이용하여 확률 구하기
% figure: tikz:fig-114-06
% answer: 
% answer_source: 
% variant_level: 0 (원본 전사)
% 이 파일은 build.mjs 가 items.json 에서 만든다. 고칠 때는 items.json 을 고친다.
\smash{\raisebox{1.5mm}{\dmkichul{베이직쎈 확률과 통계 114쪽 06번}}}
\noindent\begin{minipage}[t]{0.58\linewidth}\vspace{0pt}\raggedright 정규분포 $\mathrm{N}(m{,}\mkern3mu\mathopen{}\sigma^2)$을 따르는 확률변수 $X$에 대하여 $\mathrm{P}(m\le X\le x)$는 오른쪽 표와 같다. $X$가 정규분포 $\mathrm{N}(32{,}\mkern3mu\mathopen{}3^2)$을 따를 때, 위의 표를 이용하여 $\mathrm{P}(X\ge a)=0.0228$을 만족시키는 상수 $a$의 값을 구하면?\end{minipage}\hfill\begin{minipage}[t]{0.4\linewidth}\vspace{0pt}\centering \resizebox{\ifdim\width>\linewidth\linewidth\else\width\fi}{!}{% 114-06(인쇄 114쪽) — 114-06 오른쪽에 놓인 표($x$ / $\mathrm{P}(m\le X\le x)$).
% 스캔 표라 크롭하지 않고 tabular 로 다시 조판했다(칸 값은 원문 그대로 · 원문에 없는 행은 넣지 않는다).
{\setlength{\tabcolsep}{5pt}\renewcommand{\arraystretch}{1.25}%
\begin{tabular}{|c|c|}
\hline
$x$ & $\mathrm{P}(m\le X\le x)$ \\
\hline
$m+\sigma$ & $0.3413$ \\
\hline
$m+1.5\sigma$ & $0.4332$ \\
\hline
$m+2\sigma$ & $0.4772$ \\
\hline
\end{tabular}}
}\end{minipage}\par
\dmchoicesauto{}{$35$}{$36$}{$37$}{$38$}{$39$}
